% Geometry of the collision (Sec. II A). Oblique projection: lab x up, lab z to the right,
% lab y drawn along (-0.75,-0.3). Discs normal to z use the transform cm={0,1,-0.75,-0.3,(centre)}.
\begin{tikzpicture}[>=stealth, font=\footnotesize]
  \tikzset{disc/.style={cm={0,1,-0.75,-0.3,#1}}}
  % emitted coherent beam (drawn first, behind everything)
  \fill[violet!12] (0,1.3) -- (7.2,1.6) -- (7.2,-1.6) -- (0,-1.3) -- cycle;

  % optical axis
  \draw[->, gray] (-1.2,0) -- (8.6,0) node[right, black] {$z$};

  % electron sheet
  \filldraw[disc={(0,0)}, fill=blue!25, draw=blue!60!black] (0,0) circle (1.3);
  \node[blue!50!black, align=center] at (-0.3,1.6) {Mănunchi de electroni};
  \draw[->, very thick, blue!60!black] (0.45,1.05) -- (1.25,1.05) node[right] {$\beta$};

  % LG pulse travelling along -z, drawn as a helical wavefront under a Gaussian envelope.
  % Helix parameter t: axial position hz(t), radius hr(t); a point (z, r cos t, r sin t) is projected to
  % (z - 0.55 r sin t, r cos t - 0.2 r sin t). Lab y points towards the viewer, so sin t > 0 is in front.
  \tikzset{declare function={hz(\t)=1.9+0.066*\t; hr(\t)=0.08+0.7*exp(-((hz(\t)-3.15)/0.6)^2);}}
  \foreach \k in {0,...,5} {% back halves first, faint
    \draw[red!30, line width=1.1pt, line cap=round, domain={2*pi*\k+pi}:{2*pi*\k+2*pi}, samples=40,
    variable=\t, smooth]
    plot ({hz(\t)-0.55*hr(\t)*sin(\t r)}, {hr(\t)*cos(\t r)-0.2*hr(\t)*sin(\t r)});
  }
  \foreach \k in {0,...,5} {% front halves on top, with a white halo so they pass over the back
    \draw[white, line width=3pt, line cap=round, domain={2*pi*\k+0.15}:{2*pi*\k+pi-0.15}, samples=40,
    variable=\t, smooth]
    plot ({hz(\t)-0.55*hr(\t)*sin(\t r)}, {hr(\t)*cos(\t r)-0.2*hr(\t)*sin(\t r)});
    \draw[red!75!black, line width=1.6pt, line cap=round, domain={2*pi*\k}:{2*pi*\k+pi}, samples=40,
    variable=\t, smooth]
    plot ({hz(\t)-0.55*hr(\t)*sin(\t r)}, {hr(\t)*cos(\t r)-0.2*hr(\t)*sin(\t r)});
  }
  \draw[->, very thick, red!70!black] (3.75,-1.15) -- (2.55,-1.15);
  \node[red!60!black, align=center] at (3.15,-1.55) {Puls LG$_p^{l_L}$: $\omega_L$, $a_0$, $w_0$};

  % screen at distance Z
  \filldraw[disc={(7.2,0)}, fill=gray!12, draw=gray!70!black] (0,0) circle (1.6);
  \node[gray!40!black] at (7.2,1.85) {Ecran};
  % cartoon of the emitted field on the screen: hue = phase l*phi + alpha*rho^2 (spiral isophases, l = 2),
  % brightness = LG_0^2 amplitude (ring of radius ~0.8). Polar patches in the screen's local coordinates.
  \begin{scope}[disc={(7.2,0)}]
    \foreach \i in {0,...,11} {
      \foreach \j in {0,...,47} {
        \pgfmathsetmacro{\ra}{0.125*\i}
        \pgfmathsetmacro{\rb}{\ra+0.125}
        \pgfmathsetmacro{\pa}{7.5*\j}
        \pgfmathsetmacro{\pb}{\pa+7.5}
        \pgfmathsetmacro{\rmid}{\ra+0.0625}
        \pgfmathsetmacro{\phm}{\pa+3.75}
        \pgfmathsetmacro{\hue}{mod(2*\phm/360+1.2*\rmid*\rmid/(2*pi)+10,1)}
        \pgfmathsetmacro{\amp}{(\rmid/0.8)^2*exp(1-(\rmid/0.8)^2)}
        \pgfmathsetmacro{\bri}{0.3+0.7*\amp}
        \definecolor{patchhsb}{hsb}{\hue,0.75,\bri}\colorlet{patchcolor}[rgb]{patchhsb}
        \filldraw[fill=patchcolor, draw=patchcolor, line width=0.15pt]
        (\pa:\ra) -- (\pa:\rb) arc[start angle=\pa, end angle=\pb, radius=\rb]
        -- (\pb:\ra) arc[start angle=\pb, end angle=\pa, radius=\ra] -- cycle;
      }
    }
  \end{scope}

  % observation angles theta and phi; P = local (1.2 cos 40, 1.2 sin 40) on the screen
  \coordinate (P) at (6.506,0.635);
  \draw[dashed] (0,0) -- (P);
  \fill[white] (P) circle (0.05);
  \draw (5.0,0) arc (0:5.6:5.0);
  \node at (5.2,0.27) {$\theta$};
  \draw[white, dashed, disc={(7.2,0)}] (0,0) -- (1.5,0);
  \draw[white, disc={(7.2,0)}] (0,0) -- (1.2*0.766,1.2*0.643);
  \draw[white, ->, disc={(7.2,0)}] (0.7,0) arc (0:40:0.7);
  \node[white] at (6.95,0.85) {$\phi$};

  % distance to the screen
  \draw[<->] (0,-2.15) -- (7.2,-2.15) node[midway, fill=white, inner sep=1pt] {$Z$};

  % laboratory axes
  \begin{scope}[shift={(-1.5,-2.0)}]
    \draw[->] (0,0) -- (0,0.6) node[above] {$x$};
    \draw[->] (0,0) -- (-0.6,-0.24) node[left, inner sep=0.5pt] {$y$};
    \draw[->] (0,0) -- (0.6,0) node[right] {$z$};
  \end{scope}
\end{tikzpicture}
